\section{Konzeptionelle Grundlagen}\label{sec:grundlagen}
Die Analyse setzt zwei Klärungen voraus: was unter einem KI-gestützten Code-Generator zu verstehen ist und welche agilen Praktiken den Bezugsrahmen bilden, an dem die folgenden Kapitel die Wirkungen prüfen.

\subsection{KI-gestützte Code-Generatoren: Funktionsweise und Einsatzformen}\label{subsec:generatoren}
Als KI-gestützte \glspl{codegenerator} gelten hier Werkzeuge auf Basis großer Sprachmodelle (\aclp{LLM}, \acs{LLM})\acused{LLM}, die auf umfangreichen Mengen \glslink{opensource}{quelloffenen} Codes trainiert sind und aus dem Kontext -- Kommentaren, Funktionsnamen, umgebendem Code -- Vervollständigungen vorschlagen.\footcite[\vglf][\pagef 754]{Pearce.2022} Verbreitet sind mehrere Modellfamilien, darunter GPT-4o, Claude 3.5, Gemini 1.5, Codestral und Llama 3.\footcite[\vglf][\pagef 7300]{Kharma.2026} Zu unterscheiden sind drei Einsatzformen: die dialogische Nutzung über eine Weboberfläche, die Integration in die Entwicklungsumgebung und \glslink{agentic}{agentische Werkzeuge}, die dort selbstständig Dateien durchsuchen, ändern und Programme iterativ \glslink{debugging}{debuggen}.\footcite[\vglf][\pagef 6]{Becker.2025} Die Unterscheidung ist keine Formalie: Studien fixieren jeweils eine Form. Eine Qualitätsuntersuchung bildet dabei bewusst die kostenfreie Weboberfläche nach, wie sie besonders Berufseinsteiger nutzen.\footcite[\vglf][\pagef 67]{Swaraj.2026} Ihre Befunde gelten folglich nicht ohne Weiteres für den in Entwicklungsumgebungen integrierten Einsatz.

\subsection{Agile Softwareentwicklung als Einsatzkontext}\label{subsec:agil}
\glslink{agile}{Agile Entwicklung} liefert Software in kurzen Abständen und reflektiert regelmäßig das eigene Vorgehen.\footcite[\vglf][o.\,S.]{Beck.2001} In \gls{scrum} verdichtet sich das zu \glspl{sprint} von höchstens einem Monat, in denen die Qualität nicht abnehmen darf.\footcite[\vglf][\pagef 8]{Schwaber.2020} Als Qualitätsmaßstab dient die \gls{dod}, eine formale Beschreibung des Zustands, den ein \gls{increment} erreichen muss; konkretisiert wird sie über Kriterien wie fehlerfreie \glspl{unittest}, ersatzweise über das Vier-Augen-Prinzip.\footcites(\vglf)()[][\pagef 13]{Schwaber.2020}[][\pagef 93]{Preussig.2020} Zwei getrennte Schleifen prüfen das Erreichte: das \gls{sprintreview} das Produkt, die \gls{retrospektive} den Prozess.\footcite[\vglf][\pagef 10 f.]{Schwaber.2020} Verantwortlich ist das Team, das für das Increment einsteht und selbst entscheidet, wie Anforderungen umgesetzt werden.\footcite[\vglf][\pagef 5, \pagef 9]{Schwaber.2020} Hinzu treten Praktiken mit hohem Anteil sozialer Interaktion wie \gls{pairprogramming}\footcite[\vglf][\pagef 291]{Neumann.2026} sowie \glspl{codereview}, deren Geschwindigkeit als Prozessgröße erhoben wird.\footcite[\vglf][\pagef 37]{DORA.2024} Diese sechs Elemente bilden das Raster für die Kapitel~\ref{sec:potenziale} und~\ref{sec:risiken}.
